% CMPE 258 final report. IEEE conference format (IEEEtran), compiled with tectonic (conda env `paper`).
% Every number in this file must come from a logged run; the source of each table/figure is given in a comment.
\documentclass[conference]{IEEEtran}
\usepackage[T1]{fontenc}
\usepackage{amsmath,amssymb}
\usepackage{graphicx}
\usepackage{booktabs}
\usepackage{xcolor}
\usepackage{tikz}
\usetikzlibrary{arrows.meta,positioning,shapes.geometric,fit,calc}
\usepackage[hidelinks]{hyperref}
\usepackage{cite}

\newcommand{\videourl}{VIDEO-LINK-PLACEHOLDER} % NEEDS_AARON: replace with the YouTube/Drive link
\newcommand{\repourl}{https://github.com/AaronChandler976/curved-path-puzzle-solver}

\begin{document}
\title{Cost-Aware Query Ordering for Assembly-by-Disassembly\\with a Tight-Clearance SE(3) Planner}
\author{\IEEEauthorblockN{Aaron Chandler}
\IEEEauthorblockA{CMPE 258 Deep Learning, San Jos\'e State University\\
Code: \href{\repourl}{\texttt{\repourl}} \quad Video: \href{\videourl}{\texttt{\videourl}}}}
\maketitle

\begin{abstract}
% 150--250 words, with the key test numbers from reports/final_eval/ (filled in Phase 5).
ABSTRACT-PENDING-FINAL-EVAL
\end{abstract}

\section{Introduction}
\label{sec:intro}
INTRO-PENDING

\noindent\textbf{Contributions.}
\begin{itemize}
\item CONTRIB-PENDING
\end{itemize}

\section{Related Work}
\label{sec:related}
\textbf{Assembly-by-disassembly and learned part ordering.}
Assemble Them All~\cite{tian2022ata} plans assembly by disassembling real CAD assemblies with a physics-based
planner and released a large Fusion~360 benchmark, which our data comes from. Cebulla et al.~\cite{cebulla2023} showed
that a GNN predicting each part's removability can reorder the removal tests of AbD and cut failed attempts. ASAP~\cite{tian2024asap}
guides a disassembly tree search with a GNN (PointNet~\cite{qi2017pointnet} part features and GATv2~\cite{brody2022gatv2}
layers) trained on simulator labels, and adds gravitational stability. GRACE~\cite{atad2023} and graph-transformer
models predict sequences step by step. These methods learn \emph{whether} a part comes out and count tests. None of
them models \emph{how long} a test takes or treats a planner timeout as a censored observation, and their testers are
discretised physics simulators rather than a continuous tight-clearance planner.

\textbf{Tight-clearance motion planning.} CMG-RRT~\cite{livnat_cmgrrt} plans the separation of tightly fitting parts by
following the critical manifold of contact with an SDF-based oracle. It handles motions that need rotation, but a
query can take from milliseconds to the full budget, which is what makes query ordering matter.

\textbf{Predicting planner cost.} Odense et al.~\cite{odense2022} predict the runtime of sampling-based planners with a
GNN to pick a planner from a portfolio. Jiang et al.~\cite{jiang2026checkorder} order narrow-phase collision checks by an
expected-time criterion with learned priors, one level below our queries. We bring cost prediction to the AbD test
order. Learned survival models~\cite{gensheimer2019} give both quantities that the classical optimal order for
sequential tests needs~\cite{simon1975}: the success probability and the expected (censored) cost.

\textbf{Learning to rank.} Our tail-aware objective relaxes the order's replay cost with NeuralSort~\cite{grover2019neuralsort};
a listwise ranking loss~\cite{cao2007listnet} is an exploratory alternative.

\section{Method}
\label{sec:method}
\subsection{Problem: assembly-by-disassembly with expensive queries}
An assembly is a set of rigid parts $\mathcal{P}$ in their assembled poses. A state $S\subseteq\mathcal{P}$ is the set
of parts still present. At each state, assembly-by-disassembly (AbD)~\cite{tian2022ata} looks for one part $i\in S$
that the motion planner can move out of the others' way; the part is removed and the search continues on
$S\setminus\{i\}$, so the reversed sequence is an assembly plan. Each test (a \emph{query}) is expensive: our oracle
returns \textsc{success} with a collision-free SE(3) path or \textsc{unsolved} when its budget $T$ runs out. Because a
sampling planner cannot prove that a part is stuck, an unsolved query is a right-censored observation of the time the
planner would have needed, not an infeasibility label. The cost of a state under a query order $\pi$ is
\begin{equation}
C(\pi;S)=\sum_{k=1}^{n} t_{\pi_k}\prod_{j<k}\bigl(1-y_{\pi_j}\bigr),
\label{eq:cost}
\end{equation}
the planner time spent until the first success ($y_i\in\{0,1\}$ success, $t_i$ query time). The ranker's job is to
choose $\pi$ so that the total over all visited states is small.

\subsection{Oracle: straight-pull pre-check, then CMG-RRT}
The oracle (``oracle v2'') tests one moving part against the union of the remaining parts as a single static
obstacle. Collision is decided with an exact per-part signed distance field (mesh distance with a winding-number
sign); the union's SDF is the minimum of the parts' SDFs, and penetration below a contact tolerance
$\varepsilon=0.005$ (unit frame) counts as contact. Stage~1 sweeps the part straight along each candidate escape
direction, certified at a step of at most $\varepsilon/2$. If no straight pull is free, stage~2 runs
CMG-RRT~\cite{livnat_cmgrrt}, a tight-clearance planner that follows the contact manifold with rotations, unchanged
except for our collision adapter. Every success is re-validated edge by edge at step $\le\varepsilon/2$. The query
cost is the planner wall-clock time of both stages; the budget is $T=30$\,s.

\subsection{Ranker: GNN with a censoring-aware survival head}
Each state is a contact graph: nodes are present parts, directed edges join parts in contact. Node features (18) are
log volume, area, bounding-box extents, position relative to the sub-assembly centre and its norm, contact count and
fraction, a 6-bin contact-normal histogram and the number of present parts; edge features (5) are the mean
``away'' direction, contact fraction and contact area. The encoder is a residual GINE network~\cite{hu2020pretrain}
with LayerNorm. The head predicts, for $K=8$ log-spaced time bins $e_1<\dots<e_K=T$, the discrete hazard
$h_{ik}=\sigma(z_{ik})$ of part $i$ succeeding in bin $k$~\cite{gensheimer2019}. With $b_i$ the bin of an observed
success and $c_i$ the censoring time of an unsolved query, the loss is
\begin{equation}
\mathcal{L}_{\text{surv}}=-\sum_{i:\,y_i=1}\Bigl[\sum_{j<b_i}\log(1-h_{ij})+\log h_{ib_i}\Bigr]
-\sum_{i:\,y_i=0}\sum_{j:\,e_j\le c_i}\log(1-h_{ij}).
\end{equation}
From the survival curve $S_{ik}=\prod_{j\le k}(1-h_{ij})$ we read $p_i=1-S_{iK}$ (success within $T$) and the expected
query cost $c_i=\mathbb{E}[\min(t_i,T)]=\sum_k w_k (S_{i,k-1}+S_{ik})/2$. Parts are queried in descending $p_i/c_i$,
which minimises the expected cost of finding the first success among independent tests~\cite{simon1975}.

\subsection{Tail-aware objective}
Replay totals are dominated by a few states in which unsolved queries are tried before the first success. The
tail-aware variant adds the replay cost~(\ref{eq:cost}) of the model's own order, made differentiable with the
NeuralSort relaxation~\cite{grover2019neuralsort}: with scores $s_i=\log p_i-\log c_i$ and the soft permutation
$\hat P$, $\tilde y=\hat P y$, $\tilde t=\hat P t$ and
$\mathcal{L}=\mathcal{L}_{\text{surv}}+\lambda\,\overline{\sum_k \tilde t_k\prod_{j<k}(1-\tilde y_j)}/T$
($\tau=0.1$, $\lambda\in\{0.1,1\}$).

\subsection{Data: valid states from known feasible disassemblies}
Following the course feedback, the states are not generated by random backward rollouts. They come from
\emph{known feasible disassemblies} of real CAD products. We start from the assembled Fusion~360 model (Assemble Them
All~\cite{tian2022ata}), query \emph{every} present part at each visited state, remove one planner-verified success
chosen uniformly at random (logged seed), and repeat. Every visited state is therefore reachable by a validated
removal sequence, and every query in it has a label: its outcome and its cost under both oracle stages, so the
``pre-check $\rightarrow$ CMG-RRT'' and ``CMG-RRT only'' configurations can both be replayed.
The label set \texttt{pilot\_v2} has 2{,}022 assemblies with 3--12 parts and 41{,}913 queries ($T=30$\,s, 12.6\,h on 10
CPU workers). Three sources make it up: 1{,}328 assemblies with 3--8 parts, 255 clean assemblies with 9--12 parts,
and 439 assemblies whose CAD has small interpenetrations. For those we allow a per-pair penetration up to the
measured start depth plus $\varepsilon$, under validated caps; 919 assemblies exceed the caps and are excluded.
A machine-load guard pauses dispatch while the desktop is busy, so query times are clean timing labels. 5.1\% of
queries overlapped load and were re-run; none remained confounded.
The split groups near-duplicate assemblies (identical or near-identical part sets) and puts \emph{only fresh
assemblies}, never used in development, into the test set: 347 test assemblies (51\% with 9--12 parts, 63\% with
tolerated interpenetration), 250 val and 1{,}425 train. The test set thus measures transfer to larger and harder
products. The evaluation protocol was pre-registered before any model was trained on these data.

\section{Experiments}
\label{sec:exp}
EXP-PENDING

\section{Failure Analysis, Negative Results and Limitations}
\label{sec:fail}
FAIL-PENDING

\section{Conclusion and Future Work}
\label{sec:concl}
CONCL-PENDING

\bibliographystyle{IEEEtran}
\bibliography{refs}
\end{document}
